\section{Problem Formulation and Data Pipeline}
\label{sec:problem}

\subsection{Supervised Reconstruction Task}

For each local window, the input is raw CSI
\(\mathbf{X}\in\mathbb{C}^{A\times S\times T_x}\), with \(A=4\) receiver
streams, \(S=242\) usable subcarriers, and \(T_x=370\) packets. Real and
imaginary components are exposed as separate numeric channels, so a batch has
shape
\[
    [B,4,242,370,2].
\]
The target \(\mathbf{Y}\in[\floorval,1]^{B\times4\times340\times100}\)
contains one normalized 100-bin Doppler spectrum per antenna and output time.
A 31-packet analysis window requires 15 packets of context on either side, so
370 packets align exactly with 340 valid target frames. The objective is to
learn \(f_\theta(\mathbf{X})\approx\mathbf{Y}\) without exposing activity
labels to the student.

\subsection{Teacher and Reference Distributions}

The official precomputed targets follow the SHARP chain
\cite{meneghello2023sharp}: packet-level normalization, sparse path estimation,
selection of the strongest path as phase reference, affine phase sanitization,
Hann-windowed Doppler analysis, per-frame normalization, and flooring at
\(\floorval\approx0.0631\). We pair these paper-sourced maps with their raw CSI,
and use the same map distribution both for regression metrics and for the frozen
HAR classifier trained by the provided reproduction notebook.

The final experiment uses 92 temporally aligned recordings from all 12
available legacy AR scenario folders in the public 80-MHz CSI dataset
\cite{meneghello2023dataset}. Five of 97 candidate pairs are excluded because
their raw and Doppler lengths do not have SHARP's expected temporal offset;
forcing them into training would misalign input and target frames. Within every
remaining recording, the first \(80\%\) is available for training,
\(80\)--\(90\%\) for validation, and \(90\)--\(100\%\) for final testing. All
12 scenarios contribute to training and source validation; classifier-compatible
S1a/S1b/S1c determine early stopping and final reporting. A 31-frame guard
between intervals prevents context leakage. Training windows use stride 170,
while frozen-classifier evaluation uses stride 30 for denser coverage.
Uncompressed NumPy memory maps and deterministic recording-aware windowing make
this protocol reproducible without changing its samples.

\section{Learning Framework}
\label{sec:framework}

\subsection{Architecture Progression}

The initial temporal U-Net flattened antenna, subcarrier, and complex
components into channels, then assigned an independent output channel to every
antenna--Doppler pair. This makes a recording-average vertical spectrum an easy
solution. The next model shared one single-antenna network across all antennas
and added a spatial output head, but retained a 25-bin spectral bottleneck and
interpolation. The final model removes both shortcuts. \tabref{tab:models}
records the representative configurations and W\&B runs used in the report.

\begin{table*}[t]
\centering
\caption{Representative student architectures. Parameter counts are measured
from the saved configurations.}
\label{tab:models}
\footnotesize
\setlength{\tabcolsep}{4pt}
\begin{tabular}{@{}p{0.20\textwidth}p{0.47\textwidth}p{0.09\textwidth}p{0.16\textwidth}@{}}
\hline
Model (representative run) & Main inductive bias and limitation & Parameters & Training domain \\
\hline
Temporal 1-D U-Net (baseline) &
Frequency and antenna axes flattened into channels; independent temporal
output channels make the mean spectrum easy to represent. &
3.245M & PI, 30 carriers \\
Shared-antenna spatial head (\texttt{hpm4mjl4}) &
One antenna-independent temporal core and a 2-D head; retains a coarse
25-bin grid followed by spectral interpolation. &
0.793M & legacy AR, 242 carriers \\
Full-resolution shared 2-D (\texttt{az9k6ori}) &
Joint frequency--time encoder, GroupNorm, direct 100-bin grid, full 2-D
decoder, and exact valid-time crop. &
1.314M & legacy AR, 242 carriers \\
\hline
\end{tabular}
\end{table*}

\subsection{Full-Resolution Shared-Antenna Model}

The final architecture processes antennas independently with identical weights.
It reshapes
\[
[B,A,S,T,2]\rightarrow[BA,2,S,T],
\]
so antenna identity cannot become a shortcut, but every antenna--window pair is
retained. A \(9\times7\) frequency--time stem and two strided 2-D stages produce
feature tensors of sizes
\[
\begin{split}
[BA,64,61,370]&\rightarrow[BA,96,31,185]\\
&\rightarrow[BA,128,16,92].
\end{split}
\]
Residual blocks use GroupNorm, avoiding recording-composition-dependent
BatchNorm statistics. Only after these convolutions have modeled local
subcarrier structure is the remaining frequency axis averaged; a standard
\(1\times1\) projection then mixes latent channels.

The bottleneck is projected directly to
\([BA,16,92,100]\). A proper 2-D time--Doppler decoder upsamples time through
92, 185, and 370 positions and merges two projected encoder skips. It predicts
all 100 Doppler bins directly, with no 25-to-100 spectral interpolation. The
output head produces \([BA,370,100]\); removing exactly 15 frames from each
side yields the valid \([B,A,340,100]\) output. \figref{fig:architecture}
summarizes this flow. The model has 1,314,385 trainable parameters. These
changes correct the identified decoder and alignment flaws, but they do not
impose complex phase equivariance.

\begin{figure*}[t]
    \centering
    \includegraphics[width=\textwidth]{architecture.pdf}
    \caption{Tensor flow of the final student. Antennas share one network;
    subcarrier structure is processed before average pooling; the decoder
    predicts the complete time--Doppler grid and applies the exact 31-packet
    valid crop.}
    \label{fig:architecture}
\end{figure*}

\subsection{Motion-Aware Objective and Sampling}

Most target energy lies near zero Doppler, so plain MSE rewards a stationary
center ridge. We use
\begin{equation}
\mathcal{L}=\mathcal{L}_{\mathrm{map}}
 +0.25\mathcal{L}_{\mathrm{motion}}
 +0.25\mathcal{L}_{\mathrm{leak}}
 +0.05\mathcal{L}_{W}.
\label{eq:loss}
\end{equation}
\(\mathcal{L}_{\mathrm{map}}\) is MSE over all antennas, times, and 100 bins.
Bins 45--55 are designated stationary. Outside them,
\(\mathcal{L}_{\mathrm{motion}}\) weights squared error continuously by target
activity \(\operatorname{clip}((Y-\floorval)/(1-\floorval),0,1)\).
\(\mathcal{L}_{\mathrm{leak}}\) penalizes predicted power above the floor where
the target is at the floor. Finally, \(\mathcal{L}_{W}\) normalizes positive
off-center power into a distribution and averages the absolute difference
between predicted and target cumulative sums along the 100-bin Doppler axis.
This one-dimensional Wasserstein term distinguishes a nearby displaced peak
from equally wrong power far away.

Training uses deterministic motion stratification. Within each recording, the
top quartile of windows by active off-center frame fraction is marked
motion-rich. Every epoch includes all 1,748 rich windows and an equal-size,
non-duplicated rotating subset of ordinary windows. Batches contain 16 rich and
16 ordinary windows drawn across eight recordings. Validation and test remain
unmodified. The optimizer is Adam with learning rate \(10^{-3}\), automatic
mixed precision, a 50-epoch limit, and patience 8 on S1a/S1b/S1c validation
loss. Training completed all 50 epochs and restored the best checkpoint from
epoch 45. The smaller comparison model uses all stride-30 windows and batch
size 256.
